\documentclass[11pt]{article}

\usepackage[margin=1in]{geometry}
\usepackage[utf8]{inputenc}
\usepackage[T1]{fontenc}
\usepackage{amsmath}
\usepackage{amssymb}
\usepackage[numbers]{natbib}
\usepackage{array}
\usepackage{booktabs}
\usepackage{enumitem}
\usepackage{xcolor}
\usepackage{caption}
\usepackage[colorlinks=true,linkcolor=blue,citecolor=blue,urlcolor=blue]{hyperref}

\title{\textbf{Cross-Engine Consensus in AI-Generated Brand Recommendations:\\
An Empirical Study Across Seven Cities and Five Large Language Models}}

\author{
Constantin Daniel\\
BrandGEO\\
Santa Cruz de Tenerife, Canary Islands, Spain\\
\texttt{support@getbrandgeo.com}
}

\date{July 2026}

\begin{document}

\maketitle

\begin{abstract}
As consumers increasingly ask conversational AI systems for commercial recommendations
(\textit{``which accounting software should I use,''} \textit{``which law firm handles
employment disputes''}), a brand's visibility in these systems has become a measurable
and commercially consequential property, distinct from classical search engine ranking.
This paper reports an empirical, multi-city study of \emph{cross-engine consensus}: the
degree to which independently queried, commercially deployed AI assistants converge on
the same named brand or provider in response to the same buyer-intent question. We
collected 278 production API responses across 56 buyer-intent prompts spanning seven
cities (London, Berlin, Madrid, New York, Paris, Rome, and Dublin) and five widely used
AI systems (OpenAI ChatGPT, Google Gemini, Anthropic Claude, Perplexity, and Meta AI),
using each vendor's own production API rather than simulated or scraped output. One
engine (ChatGPT) is excluded from analysis after an account-level API quota failure
affected all 56 of its scheduled calls; the remaining 222 responses (four engines) form
the analytic dataset. We find that categories with a small, well-documented, easily
verifiable set of providers (project management software, business banking applications,
large corporate law firms) converge strongly, frequently to unanimous agreement, while
categories populated by many similarly qualified, less differentiated local or
individual providers (employment law, independent financial advice) rarely converge at
all. Two further findings extend this picture. First, query language can change not
merely the ranking but the entire considered set of competitors: a wealth-management
query in Paris returns boutique French firms when asked in French and major
international private banks when asked in English, with almost no overlap between the
two sets. Second, one city (Rome) shows zero categories reaching even three-engine
agreement, the only such case in the dataset, while simultaneously showing the strongest
single-engine, cross-lingual stability observed anywhere in the study (one engine
reproduced an identical, identically ordered five-item ranking in both Italian and
English). We report these findings with their accompanying data-quality caveats in
full, and discuss their implications for brand visibility measurement and for the
emerging practice of Generative Engine Optimization.
\end{abstract}

\noindent\textbf{Keywords:} generative engine optimization, AI visibility, large
language models, brand recommendation, retrieval-augmented generation, cross-engine
consensus, conversational search

\section{Introduction}

Consumers are increasingly using conversational AI systems, rather than a traditional
search engine results page, as the first point of contact when researching a commercial
purchase decision. Industry measurement of this shift is still young but directionally
consistent: independent analyses report that a meaningful and growing share of
commercial-intent queries now trigger an AI-generated answer rather than a list of blue
links, and that the resulting answers frequently name specific brands or providers
without the user ever visiting a ranked list of websites \citep{semrush2026,
brightedge2026}. This shift has produced a new discipline, commonly referred to as
Generative Engine Optimization (GEO), concerned with whether and how a brand is
mentioned, ranked, or recommended inside an AI system's generated response, as distinct
from its position in a classical search index \citep{aggarwal2024geo}.

A property of this new environment that has received comparatively little direct,
side-by-side empirical attention is \emph{cross-engine consensus}: whether different AI
systems, queried independently with the same buyer-intent question, tend to agree on
which brand or provider to recommend. This question matters for at least three reasons.
First, a brand that appears consistently across engines has a materially different risk
profile than one that appears well in a single engine and is invisible everywhere else;
a monitoring or optimization strategy built around only one engine could produce a
badly distorted picture of a brand's real standing. Second, consensus (or its absence)
is itself informative about the underlying mechanism: categories where independent
systems converge look like cases where a genuinely dominant, well-documented answer
exists in the training and retrieval data available to every engine, while categories
where systems disagree look more like cases where each engine is making an
under-determined choice among many similarly qualified options. Third, for a company
attempting to measure or improve its own AI visibility, knowing whether a category is
structurally prone to consensus or fragmentation changes what a realistic goal even
looks like.

This paper reports a pilot-scale, empirical measurement of cross-engine consensus using
a dataset collected through BrandGEO's own production monitoring pipeline: a system
that fires real, commercially phrased buyer-intent prompts against the production APIs
of five widely used AI assistants and records the resulting responses. We use this
pipeline's own operational data, collected across seven cities as part of an internal
research initiative, as the empirical basis for this study, and we report both the
substantive findings and the pipeline's known data-quality limitations in full, rather
than presenting only a cleaned final result.

The remainder of this paper is organized as follows. Section~2 situates this work
relative to prior research on generative engine optimization and AI-driven search
behavior. Section~3 describes the dataset, collection methodology, and the operational
definition of cross-engine consensus used throughout. Section~4 presents results,
first in aggregate and then per city, including two extended case findings (language-
dependent competitor sets in Paris, and the zero-consensus/high-stability anomaly in
Rome). Section~5 discusses implications for brand visibility measurement and GEO
practice. Section~6 states limitations, ethical considerations, and a reproducibility
statement. Section~7 concludes.

\section{Related Work}

The term \emph{Generative Engine Optimization} was introduced by Aggarwal et al. \citep{aggarwal2024geo},
who formalized the problem of improving a piece of content's visibility inside
generative engine responses, introduced a benchmark (GEO-bench) of realistic user
queries paired with the sources needed to answer them, and showed that specific
content-level interventions (adding citations, statistics, or authoritative quotations)
could measurably increase a source's visibility in generated answers, with effects
varying substantially by domain. That work established the underlying premise this
paper builds on: visibility inside a generated answer is a distinct, measurable,
optimizable property, not a proxy for classical search rank. Our contribution differs
in unit of analysis: rather than measuring whether a piece of content is cited, we
measure whether independently queried commercial AI assistants agree with one another
on which named brand or provider to recommend for a given buyer-intent question, across
a geographically and linguistically varied set of prompts.

A second relevant strand of work concerns how retrieval-augmented and search-grounded
language models select and surface sources at all. Vendor documentation for the systems
used in this study describes materially different retrieval mechanisms: Google's Gemini
uses a search-grounding step gated by an internal confidence classifier that decides,
per query, whether to invoke a live web search or answer from parametric knowledge
alone \citep{googlegrounding2026}; OpenAI's web-enabled ChatGPT models invoke a
dedicated web search tool that can additionally be scoped to a user's approximate
location \citep{openaisearch2026}; Perplexity is built around retrieval and citation by
design. These architectural differences are a plausible mechanistic explanation for
why the same buyer-intent question can produce different answer sets across engines,
and they motivate treating engine identity, not just query content, as a variable of
interest.

Independent industry measurement of AI-driven commercial search has grown quickly
alongside these systems, though it remains largely proprietary and organization-specific
rather than peer-reviewed. Reported findings in this literature, cited here as
context rather than as directly comparable measurements, include analyses showing that
a large share of AI-cited brands correlate more strongly with unstructured web mentions
than with classical backlink counts \citep{ahrefs2026correlation}, that a majority of
locally oriented AI answers substantially under-recommend real local businesses
relative to their standing in classical local search \citep{soci2026}, and that citation
freshness (how recently a cited source was published or updated) correlates with a
source's likelihood of being surfaced by a retrieval-grounded engine
\citep{amsive2026, ahrefs2026freshness}. Where this paper draws on such
industry-reported figures for context, they are cited to BrandGEO's own previously
published, sourced synthesis of the underlying reports, so that the provenance chain
from primary report to citation here is transparent and independently checkable.

\section{Data and Methods}

\subsection{Study design}

Between 9 and 10 July 2026, BrandGEO's monitoring pipeline was used to run a structured,
multi-city measurement exercise. Eight buyer-intent prompts were authored for each of
seven cities (London, Berlin, Madrid, New York, Paris, Rome, and Dublin), for a total
of 56 prompts. Prompts were written to resemble realistic, high commercial intent
questions a real consumer or business buyer might put to a conversational AI assistant
(for example, \textit{``What is the best CRM software for a small business in London?''}
or \textit{``Which law firm should I use for an employment dispute in Dublin?''}), each
targeting a specific commercial category (software, professional services, hospitality,
financial services, real estate, or food service, depending on the city).

Two cities (London and Dublin) are English-first markets and used eight distinct
English-language prompts, one per category. The remaining five cities (Berlin, Madrid,
Paris, Rome) used a paired bilingual design: four commercial categories, each phrased
once in the local language (German, Spanish, French, or Italian respectively) and once
in English, for eight prompts per city. This design was chosen specifically to allow
language to be isolated as a variable within a city, rather than confounded with the
city itself.

Every prompt was submitted, independently and without conversational context carried
between prompts, to each of five production AI systems: OpenAI's ChatGPT, Google's
Gemini, Anthropic's Claude, Perplexity, and Meta AI. Requests were geo-contextualized
(the system prompt informed each model of the relevant city or country context) to
approximate how a real user physically located in that market would be answered.

\subsection{Collection pipeline and known engine-level failures}

Responses were collected through BrandGEO's own production Netlify-hosted collection
functions, which call each vendor's real, versioned API and log every attempt, including
failures, to a structured results table. This is a deliberate design choice relevant to
this paper's reliability: every one of the 56 prompts targeted five engines, for 280
possible engine-level responses; 278 rows were actually recorded (two rows, both for
Meta AI in New York and Paris respectively, were lost to a transient collection gap
rather than a substantive non-response).

Of the 278 recorded rows, 56 (exactly one for each of the 56 prompts, entirely
concentrated in one engine, ChatGPT) carry a recorded \texttt{error} status rather than
a completed response. Root-cause investigation (documented in the pipeline's internal
engineering log) attributes this to a sustained, account-level API quota exhaustion on
the OpenAI account used for collection, spanning the entire window in which this study's
data was gathered. This is a genuine data-collection failure, unrelated to any
substantive property of ChatGPT's actual visibility behavior, and we exclude ChatGPT
from all analysis in this paper rather than treat its absence as a null result. We
report this openly, rather than silently omitting the engine, because a reader
attempting to reproduce or extend this work needs to know that a five-engine design
produced a four-engine analytic dataset for a specific, documented, and since-resolved
operational reason, not because ChatGPT was deliberately excluded from the study design.

The analytic dataset therefore consists of 222 completed responses across four engines
(Gemini, Claude, Perplexity, Meta AI), spanning 56 prompts in seven cities.
Table~\ref{tab:dataset} summarizes engine-level completion counts.

\begin{table}[h]
\centering
\caption{Recorded engine-level responses per city. ChatGPT is excluded from analysis
(see Section 3.2); its column reports recorded error rows only.}
\label{tab:dataset}
\begin{tabular}{lccccc}
\toprule
City & Prompts & Gemini & Claude & Perplexity & Meta AI \\
\midrule
London & 8 & 8 & 8 & 8 & 8 \\
Berlin & 8 & 8 & 8 & 8 & 8 \\
Madrid & 8 & 8 & 8 & 8 & 8 \\
New York & 8 & 8 & 8 & 8 & 7 \\
Paris & 8 & 8 & 8 & 8 & 7 \\
Rome & 8 & 8 & 8 & 8 & 8 \\
Dublin & 8 & 8 & 8 & 8 & 8 \\
\midrule
\textbf{Total} & \textbf{56} & \textbf{56} & \textbf{56} & \textbf{56} & \textbf{54} \\
\bottomrule
\end{tabular}
\end{table}

\subsection{Response analysis}

Each raw response was parsed by an automated pipeline to extract, where present, a
ranked or otherwise prominently featured list of named brands or providers. This
pipeline evolved over the course of the underlying project to reduce false-positive
extractions (for example, rejecting section headings, evaluation-criteria labels, and
generic category nouns that could otherwise be mistaken for brand names). Even with
these safeguards, raw extraction is not perfectly reliable, and every city's dataset was
subject to a manual review pass in which an analyst read the underlying raw response
text directly and corrected or flagged extraction artifacts before any qualitative
finding was reported. Known, disclosed categories of correction include: rejecting
brand-like fragments that are not real companies (for example, section headings or
sentence fragments mistakenly parsed as a ranked entry); flagging entities that could
not be independently verified as currently operating (for example, one Madrid response
named a restaurant that appears to correspond to a now-closed establishment); and
flagging at least one clear geographic mismatch (a Madrid dining query returning a
restaurant located in a different city). We report findings only where the underlying
named entity survived this manual verification step, and we disclose, per city, where a
response was judged too noise-corrupted to support a reliable qualitative reading.

Gemini's responses showed a further, engine-specific pattern: on a nontrivial minority
of prompts, Gemini's raw text clearly named specific brands in prose, but the automated
extraction step returned no structured list at all, most often because Gemini answered
in a bold-text or bulleted-paragraph format that the extraction pipeline at the time
did not fully parse as a ranked list. This produced an elevated rate of
\emph{apparent} non-response for Gemini in several cities (most severely in Rome, where
five of eight prompts showed this gap, and New York, where four of eight did). We treat
this explicitly as a pipeline extraction limitation, not evidence that Gemini lacked an
opinion, and we flag it separately from genuine cross-engine disagreement wherever it
materially affects a result below.

\subsection{Operational definition of cross-engine consensus}

For each prompt, we define \emph{cross-engine consensus} as having been reached when at
least three of the engines that returned a usable, verifiable response for that prompt
independently named the same brand or provider as their top-ranked or single most
prominently featured recommendation. Where only three engines returned usable data for
a given prompt (most often because of the Gemini extraction gap described above), we
require unanimous three-of-three agreement rather than relaxing the threshold. This is
a coarse, binary construct, deliberately chosen for transparency over a continuous
similarity metric: it is directly interpretable (\textit{``did the engines agree or
not''}), at the cost of not distinguishing, for example, a narrow three-of-four
majority from a full four-of-four sweep. We report both, in text, wherever the
distinction is substantively meaningful.

\subsection{Ethics and scope}

All prompts describe realistic, commercially neutral information-seeking scenarios (for
example, choosing a law firm, an accounting package, or a hotel). No prompt was
designed to elicit content harmful to any named brand, and no named brand was
deliberately advantaged or disadvantaged in prompt design. Named brands, providers, and
professionals appearing in this paper's results are reported exactly as they appeared in
a real, production AI system's response to a realistic query; their inclusion here is
descriptive of observed AI behavior, not an endorsement, criticism, or claim about the
named entity's actual quality.

\section{Results}

\subsection{Aggregate pattern: category structure predicts consensus}

Table~\ref{tab:results} summarizes, per city, the number of categories (of eight
prompts, four categories in bilingual cities and eight in monolingual cities) that
reached cross-engine consensus under the definition in Section~3.4, alongside the
single most notable qualitative finding for that city.

\begin{table}[h]
\centering
\caption{Per-city consensus summary and headline qualitative finding.}
\label{tab:results}
\small
\begin{tabular}{>{\raggedright\arraybackslash}p{1.5cm}>{\raggedright\arraybackslash}p{2.3cm}>{\raggedright\arraybackslash}p{4.4cm}>{\raggedright\arraybackslash}p{5.4cm}}
\toprule
City & Consensus reached & Example unanimous or near-unanimous pick & Headline finding \\
\midrule
London & 6 of 8 categories & Capsule CRM (4/4); Starling, Monzo, Revolut (each 4/4) & Small, well-documented software and fintech categories converge strongly; employment law and independent financial advice do not. \\
Berlin & 2 of 4 categories (bilingual) & Lexware/Lexoffice (4/4, German); Personio (3/4, German) & Gemini and Meta AI both returned zero competitors in English for a prompt they answered fully in German: a language-dependent visibility gap, not just re-ranking. \\
Madrid & 1 strong, 1 partial of 4 categories & Hilton Madrid Airport (unanimous, both languages) & The strongest single consensus item in the dataset; a bidirectional, engine-specific language gap: Gemini silent in Spanish/complete in English, Meta AI the exact opposite, same prompt. \\
New York & Majority of 8 categories & StreetEasy (3/4, unanimous among responders) & Institutions (Douglas Elliman, Corcoran) beat individually named experts; national law firms beat city-specific boutiques for startup incorporation. \\
Paris & 1 of 4 categories, unanimous & Qonto (4/4, both languages, zero exceptions) & Query language changes the entire considered competitor \emph{set}, not just the ranking, for wealth management: French queries surface boutique French firms, English queries surface major international private banks. \\
Rome & 0 of 4 categories & (none reached three-engine consensus) & The only city with zero cross-engine consensus in any category, yet the strongest single-engine, cross-lingual stability in the study (Meta AI reproduced an identical ranked list in two languages). \\
Dublin & 5 of 8 categories, 2 fully unanimous & Monday.com (4/4); Revolut Business (4/4) & Highest full-consensus rate in the study, alongside a distinct engine-level failure: Meta AI showed a recurring bias toward large, well-known institutions regardless of whether the query asked for a boutique or independent option. \\
\bottomrule
\end{tabular}
\end{table}

The clearest structural pattern across all seven cities is that categories with a small
number of well-documented, easily distinguished, and widely reviewed providers (project
management software, business banking, established corporate law firms) converge, often
to full unanimity, while categories populated by many similarly qualified, less
differentiated providers, particularly individually named professionals (employment
solicitors, independent financial advisers, personal injury attorneys), rarely converge
at all. This pattern held independently in London, New York, and Dublin, three cities
that did not share a bilingual design and therefore cannot be explained by a language
confound.

\subsection{Institutions over individuals: New York}

New York's dataset adds a specific variant of this pattern. For real estate brokerage,
three long-established, large brokerages (Douglas Elliman, Corcoran, Brown Harris
Stevens) each achieved three-of-four engine consensus, while individually named
``star'' agents and boutique teams named in the same responses never achieved
cross-engine agreement above one engine. A parallel pattern held for legal services:
for startup incorporation, two large, nationally recognized technology-sector law firms
(Fenwick \& West, Cooley) reached three-of-four consensus ahead of any New York
specific boutique firm, extending a national-provider-over-local-provider pattern
already observed in this project's prior industry-level content research. Gemini's
extraction gap was most pronounced here among the monolingual cities, returning no
structured competitor list on four of eight prompts despite naming brands in prose,
underscoring the pipeline limitation flagged in Section~3.3.

\subsection{Language changes the competitor set, not just the order: Paris}

Paris produced the strongest single unanimous result in the entire dataset: for the
prompt asking which online bank to use as a French small business, every engine that
returned a usable response named Qonto as its top pick, in both the French-language and
English-language runs, with no exceptions across four engines and two languages.

The city's most theoretically significant finding, however, concerns a different
category: wealth management. Here, query language did not merely reorder a fixed set of
candidate firms, it changed which firms were considered at all. French-language queries
converged on independent, boutique French wealth-management firms (most notably Cheval
Blanc Patrimoine, named in three of four French-language responses and in zero
English-language responses). English-language queries, by contrast, converged on major
international private banking divisions (Soci\'et\'e G\'en\'erale Private Banking, BNP
Paribas Wealth Management), which were near-unanimous once both language runs were
combined but almost entirely absent from the French-language responses. This is a
substantively different phenomenon from the presence/absence gaps observed in Berlin
and Madrid: rather than one engine going silent in one language, every engine that
answered in a given language converged on a systematically different category of
provider depending on that language alone.

\subsection{The Rome anomaly: zero consensus, maximal within-engine stability}

Rome is the only city in the dataset in which no category reached even three-engine
agreement. Every one of its four categories (boutique hotels, business-dinner
restaurants, real estate agencies, and a hotel near Roma Termini station) fragmented
across engines. Rome also showed the most severe Gemini extraction gap in the study
(five of eight prompts returned no structured competitor list), and one Perplexity
response was judged, on manual review, to be the most heavily noise-corrupted single
response encountered anywhere in this research program (of eight extracted entries,
only one was a genuine restaurant name; the remainder were street addresses, generic
labels, and a reservation platform mistakenly parsed as a competitor).

Despite this absence of cross-engine agreement, Rome showed the single strongest
within-engine, cross-lingual stability result in the study. For the business-dinner
restaurant category, Meta AI produced a five-item ranked list in Italian and, when asked
the equivalent question in English, produced the identical five restaurants in the
identical order: a full, exact reproduction rather than merely a similar or overlapping
set. Perplexity showed a comparable, slightly weaker pattern for real estate agencies,
naming essentially the same nine agencies in both languages with only reordering
between them. No other engine showed this pattern in Rome, and no engine showed a
comparably exact reproduction in any other city.

\subsection{An engine-specific generalization failure: Dublin}

Dublin, an English-only monolingual city, showed both the strongest overall consensus
rate in the study (a corporate legal dispute prompt converged on Ireland's five largest
commercial law firms, all named by all three responding engines; a project management
software prompt and a business banking prompt each reached full four-of-four unanimity,
on Monday.com and Revolut Business respectively) and a distinct, engine-specific
failure mode not observed elsewhere. Meta AI's answers to three of eight prompts,
covering financial advisers, wealth management for expatriates, and property-purchase
solicitors, were each nearly identical in composition to its answer for an unrelated
prompt about large-scale commercial litigation, naming the same set of large,
well-known institutions (for example, Davy, Investec, Cantor Fitzgerald, and Goodbody
Stockbrokers for two nominally distinct wealth-related queries) regardless of whether
the query explicitly asked for an independent or boutique option. Claude, Gemini, and
Perplexity, queried with the identical prompts, differentiated cleanly between these
scenarios. We interpret this as a query-intent generalization failure specific to one
engine on this class of prompt, rather than evidence about the real market structure of
Dublin's financial advisory sector, and note it as a further reason a single-engine
visibility measurement can misrepresent a brand's actual standing.

\section{Discussion}

\subsection{Implications for brand visibility measurement}

The central practical implication of these results is that a brand visibility
measurement built on a single AI engine, however prominent that engine may be, risks
systematically misrepresenting a brand's real standing. A brand that is the unanimous
top pick across four independent engines (as Qonto was in Paris, or Monday.com in
Dublin) occupies a materially more defensible market position than a brand that is
top-ranked in one engine and simply absent, rather than merely lower-ranked, in others,
a distinction that a single-engine measurement cannot express at all. Equally, the
Berlin and Madrid language-gap findings show that even a multi-engine measurement can be
misleading if it is monolingual: a brand's visibility can differ sharply by the
language a real customer happens to use, in ways that are not predictable from
visibility in any one language alone.

\subsection{Implications for category-level strategy}

The consensus-versus-fragmentation pattern observed across London, New York, and Dublin
suggests that the realistic goal for a GEO program differs meaningfully by category
structure. In categories that already show strong cross-engine consensus around a small
set of incumbents (well-documented software, large banks, large corporate law firms),
the practical opportunity for a new or smaller entrant is likely to be difficult:
displacing an already-converged consensus requires the same widely available, easily
verifiable signal that produced the consensus in the first place. In categories that
show consistent fragmentation (individually named professionals, small local providers
without a single obviously dominant option), the absence of a converged answer suggests
more genuine opportunity for a well-documented, easily verifiable provider to become the
engines' converged answer, precisely because no such consensus yet exists to displace.

\subsection{Threats to validity}

This study has several limitations that bound the strength of any conclusion drawn from
it. It is a single-snapshot measurement, collected over a two-day window, not a
longitudinal study; AI system outputs are known to vary over time as underlying models,
retrieval indexes, and grounding configurations are updated, and a repeated measurement
at a later date could show different results for the same prompts. The sample is
small by conventional statistical standards (56 prompts, 222 analyzed responses), and
per-prompt engine counts are frequently three or four rather than five, both because of
the ChatGPT collection failure (Section 3.2) and the Gemini extraction gap
(Section 3.3); we have treated both as data-quality issues to be disclosed and worked
around, not as substantive findings, but a reader should weigh conclusions accordingly.
The manual verification step described in Section 3.3, while conducted with the
specific goal of reducing extraction artifacts, necessarily involves analyst judgment
and is not independently replicated by a second coder in this study. Finally, this
study measures a fixed, English-plus-one-additional-language design per city; it cannot
speak to consensus behavior in language pairs, categories, or geographies outside the
seven cities and eight categories actually tested.

\section{Limitations, Ethics, and Reproducibility}

We restate, for a reader relying on this paper in isolation, the two data-collection
issues disclosed in Section~3: ChatGPT is excluded from all analysis due to a
documented, account-level API quota failure affecting all 56 of its scheduled calls
during the collection window, and Gemini shows an elevated, format-dependent
non-extraction rate in several cities that should not be read as a substantive absence
of visibility. Every qualitative finding in Section~4 is drawn from responses that
survived a manual review of the underlying raw model output, and every flagged,
uncertain, or noise-corrupted data point encountered during that review is disclosed
in this paper rather than silently omitted or silently corrected. No named brand,
professional, or provider in this dataset was contacted, sampled, or otherwise involved
in the study; all data consists of AI system outputs to realistic informational
prompts, collected through each vendor's public, documented API in the ordinary course
of a commercial monitoring product's operation. The underlying per-city working notes,
including the full prompt list and every disclosed data-quality flag referenced in
Section~3.3, are maintained by the authors and available on reasonable request.

\section{Conclusion}

Across seven cities and 56 buyer-intent prompts, we find that whether independently
queried AI assistants agree on a recommended brand depends strongly on the structure of
the underlying commercial category: well-documented, small-provider-set categories
converge, often unanimously, while categories with many similarly qualified
individually named providers rarely do. We further find that query language can, in at
least one documented case, change the entire set of brands an engine considers rather
than merely reordering a fixed set, and that at least one city in our sample shows a
combination of zero cross-engine consensus with the strongest observed within-engine,
cross-lingual stability, a combination that a single aggregate consensus metric would
have obscured. These findings support a practical conclusion for both brand owners and
platform designers: meaningful AI visibility measurement requires sampling across
multiple engines and, where a market is multilingual, across multiple languages,
because both engine identity and query language are independent sources of variation
in what an AI system recommends, not noise around a single underlying truth.

\bibliographystyle{plain}
\begin{thebibliography}{9}

\bibitem{aggarwal2024geo}
P. Aggarwal, V. Murahari, T. Rajpurohit, A. Kalyan, K. Narasimhan, and A. Deshpande.
GEO: Generative Engine Optimization. \textit{Proceedings of the 30th ACM SIGKDD
Conference on Knowledge Discovery and Data Mining (KDD '24)}, 2024.
\url{https://arxiv.org/abs/2311.09735}

\bibitem{googlegrounding2026}
Google. Grounding with Google Search. Gemini API documentation, accessed July 2026.

\bibitem{openaisearch2026}
OpenAI. Web search tool documentation, accessed July 2026.

\bibitem{semrush2026}
Semrush. AI Overview and AI-driven search prevalence tracking, 2026. As reported in
BrandGEO Research, ``Schema Markup for AI Search: The Complete Implementation Guide,''
\url{https://getbrandgeo.com/bg-010.html}

\bibitem{brightedge2026}
BrightEdge. AI Overview prevalence tracking study, 2026. As reported in BrandGEO
Research, \url{https://getbrandgeo.com/bg-005.html}

\bibitem{ahrefs2026correlation}
Ahrefs. Brand mention and backlink correlation study across 75{,}000 brands, 2026. As
reported in BrandGEO Research, ``The Anatomy of an AI-Cited Brand,''
\url{https://getbrandgeo.com/bg-007.html}

\bibitem{ahrefs2026freshness}
Ahrefs. Content freshness and AI citation rate study, 2026. As reported in BrandGEO
Research, \url{https://getbrandgeo.com/bg-013.html}

\bibitem{amsive2026}
Amsive. Content freshness and AI citation research, 2026. As reported in BrandGEO
Research, \url{https://getbrandgeo.com/bg-013.html}

\bibitem{soci2026}
SOCi. Local Visibility Index, 2026. As reported in BrandGEO Research, ``AI Visibility
by Industry: Who Wins, Who Loses,'' \url{https://getbrandgeo.com/bg-008.html}

\end{thebibliography}

\end{document}
